\section{Introduction}

Trace organic contaminants -- pharmaceuticals, pesticides, biocides, UV filters
-- are ubiquitous in surface
waters~\cite{schwarzenbach2006challenge,loos2009eu}. A global river survey found
pharmaceutical contamination on every continent, exceeding safe concentrations
at a quarter of the sites~\cite{wilkinson2022pharma}. Pesticide concentrations
European regulation treats as acceptable are associated with losses of regional
stream invertebrate biodiversity~\cite{beketov2013biodiversity}; the substances
outpace the frameworks built to manage them~\cite{geissen2015emerging}.

These substances are regulated in the European Union through environmental
quality standards (EQS) under the Water Framework Directive~\cite{eu2000wfd},
transposed in T\"urkiye as \emph{\c{C}evresel Kalite Standartlar{\i}} (\c{C}KS).
Compliance is assessed by comparing measured concentrations against those
standards, a comparison less well defined than it appears.

Most such measurements are non-detections. The analyte is sought and not found.
That establishes an \emph{upper bound}, not absence: the substance may be
present below the method's limit of detection (LOD). The statement is
open-world.

Reporting practice replaces that statement with a number, usually zero or half
the limit of quantification (LOQ). Helsel showed that substitution fabricates
data and distorts every derived statistic~\cite{helsel2006fabricating}; the
point is not contested. It remains standard because \textbf{the data models in
use cannot hold the statement}. A relational table is closed-world: an absent
assertion is false, and a column of concentrations has no cell for ``present or
absent, below \num{0.02}\,\si{\micro\gram\per\litre}''. Substituting a zero is
not a rounding convention but a type error: it destroys the modality of the
measurement, not its magnitude.

Regulatory assessment implicitly requires three verdicts -- compliant,
exceeding, and \emph{cannot be determined}. The systems producing it can express
only the first two. Where the LOQ lies above the EQS, a laboratory can correctly
report ``not detected'' for a sample that may exceed the legal limit by two
orders of magnitude, and a closed-world pipeline records it as compliance.

Article~3(3b) of the consolidated Directive requires that a below-quantification
result whose limit exceeds the standard ``shall not be considered'' for status
assessment, and Article~4(1) of Directive 2009/90/EC requires that limit to sit
at or below \num{30}\,\% of the standard. Both mandate an outcome a two-valued
data model cannot represent, so neither can be audited from the data as
recorded.

We measure the cost in the European Environment Agency's Waterbase, the
monitoring Member States report under the Directive~\cite{eea_waterbase}:
\num{4190833} river station-years over \num{637} substances. Waterbase records a
below-quantification flag and the limit in force for each station, substance and
year, so where these are missing that absence is itself the measurement. We
quantify the share of samples below the quantification limit, the share of
station-years carrying neither flag nor limit, and the share assessed with a
method failing the legal performance criterion. Splitting the record at
\num{2015} separates a record-keeping failure the reporting authorities have
since repaired from an analytical one they have not; reading the disaggregated
release extends the audit to the standard Annex~I defines against a single
sample rather than against an annual mean.

The gap follows from a measurement model. SOSA/SSN and its descendants were
designed for in-situ \emph{sensing}, in which a detection limit is a property of
the device. Trace organic monitoring is \emph{sampling followed by laboratory
analysis}, in which the limit belongs to the method that produced the result
and its consequence belongs to the result itself. We present CENSO, an ontology
that carries the limit onto the result, records why a result is
an interval, treats a threshold as a conditional judgement with preconditions,
and makes compliance three-valued --- met, not met, or not determinable ---
subtyped by the reason it is not. The three are the outcomes the Directive
itself admits; the subtypes are what tell an unmeasurable standard apart from an
unreported bound. We also assess \num{21} ontologies, parsing each file rather
than reading its documentation, and the set reaches outside the water domain to
the analytical-chemistry and statistics vocabularies where these concepts would
live if they lived anywhere.

The ontology is the instrument, not the result. Without a vocabulary that
separates \emph{compliant} from \emph{cannot be determined}, none of the
quantities reported here can be counted at all.
